\chapter{LLM Prompt Injection for the Experimental Group}
\label{app:jla_prompt}

\textit{The following parameters were dynamically injected into the LLM system prompt for learners in the experimental condition based on their individual BKT profiles.}

\paragraph{Personalised Vocabulary Goals}
The student is currently learning these words relevant to this topic: 

\vspace{0.5em}
\noindent
\texttt{\{personalised\_vocabulary\}}
\vspace{0.5em}

\textbf{Instructions:}
\begin{enumerate}[noitemsep]
\item \textbf{Prioritise:} If the context allows, encourage the student to use these specific headwords and their derivatives / inflected forms.
\item \textbf{Challenge:} If the student writes a simple sentence, suggest rewriting it using one of these words. \textit{Example:} If target is 'exhausted' and they wrote 'I was very tired', say: ``Try using \textbf{'exhausted'} instead of 'very tired'.''
\item \textbf{Praise:} If they use one of these words, give specific praise.
\end{enumerate}

\paragraph{Personalised Grammar Goals}
The student is currently practising the following grammar structures. Use the provided definitions to guide your feedback.

\vspace{0.5em}
\noindent
\texttt{\{personalised\_grammar\}}
\vspace{0.5em}

\textbf{Instructions:}
\begin{enumerate}[noitemsep]
\item \textbf{Identify:} Check if the student has attempted these structures.
\item \textbf{Correct:} If they attempted but failed (e.g., missed the question mark in ``What about\dots''), use the \textbf{Rule} above to explain the mistake.
\item \textbf{Encourage:} If the student hasn't used these structures yet, prompt them naturally to do so. For example, ask a question that elicits one of these grammar patterns.
\item \textbf{Praise:} If they use one of these structures correctly, give specific praise.
\end{enumerate}

